% ============================================================================
%  Resolución de las observaciones del Informe 1 (FEP01, Artículo 46, p. 28)
%  Contiene las 100 filas de la matriz, asociadas a los subdocumentos indicados.
%  Las filas 01-15, 37-52 y 81-87 se incorporaron desde el borrador y requieren
%  cotejo con la matriz fuente, los responsables y la evidencia oficial.
% ============================================================================

\begin{tablaObservaciones}{Informe 1}

\grupoObservaciones{Aspectos generales, formalidad y cumplimiento de instrucciones}

\observacion[tipo=acepta, sub=0]{01}{Formulario T-12 ausente pese a que el Subdocumento 3 pág. 6 indica que acompaña la oferta técnica.}{Existe el archivo fuente AUDIT-Formulario-T-12.md con 42 requisitos; el PDF independiente y su incorporación a la oferta final no están disponibles en este checkout.}{AUDIT-Formulario-T-12.md; PDF final y ubicación en la oferta pendientes.}

\observacion[tipo=acepta, sub=0]{02}{Formulario T-19 diferido a la propuesta final cuando el ítem 13 de las bases lo incluye en esta instancia.}{Solo se localizó la ficha Tipo 5; no están las fichas Tipo 1 a 4 ni el formulario T-19 consolidado. La respuesta de entrega completa queda pendiente.}{formulario_t19_tipo5_innovacion.md; T-19 Tipo 1 a 4 y consolidado pendientes.}

\observacion[tipo=acepta, sub=0]{03}{El Anexo A del Subdocumento 4 no es el Formulario T-11, es una tabla de emplazamiento sin cantidades ni configuración.}{Existe el fuente AUDIT-Formulario-T-11.md. Falta cotejar sus cantidades, configuración y versión de entrega; no está el PDF final.}{AUDIT-Formulario-T-11.md; cotejo técnico y PDF final pendientes.}

\observacion[tipo=acepta, sub=0]{04}{El interior de la propuesta corresponde a una plantilla académica genérica sin identidad corporativa de empresa.}{La plantilla audIT disponible contiene fuentes parciales de S1, S2 y S9, pero no se encontró el archivo principal ni los estilos globales; no se puede verificar la aplicación integral.}{recursos/Formato-Oferta-audIT; archivo principal, estilos y compilación pendientes.}

\observacion[tipo=acepta, sub=0]{05}{Se cita a una unidad académica como autor bibliográfico en los Subdocumentos 1, 2, 3 y 13.}{Se creó referencias/referencias.bib con entradas compartidas para las citas de S1/S2 y sus fuentes normativas. El cotejo transversal de autorías/citas de S3 y S13 sigue pendiente; también falta integrar el archivo en el paquete de compilación común.}{S1 §Referencias, S2 §2.2.2; bibliografía disponible; cotejo S3/S13 e integración de plantilla pendientes.}

\observacion[tipo=acepta, sub=0]{06}{Numeración de títulos inconsistente; el índice del Subdocumento 1 muestra 1.3, 1.2.1 y 2 1.3.}{Los títulos actuales de S1 siguen la jerarquía 1.1 a 1.6; la revisión del índice y de la versión compilada queda pendiente.}{S1 §§1.1–1.6; índice y PDF final pendientes.}

\observacion[tipo=acepta, sub=0]{07}{Títulos de Nivel 1 a mitad de página con el número en cuerpo tipográfico gigante.}{La normalización de estilos no puede verificarse visualmente: no hay PDF compilado de la oferta en el checkout.}{PDF compilado de la oferta; no disponible.}

\observacion[tipo=acepta, sub=0]{08}{Título seguido inmediatamente de otro título y una tabla, sin párrafo de contexto ni conclusión.}{S1 y S2 incluyen párrafos de contexto en sus fuentes actuales; la revisión de todas las secciones, tablas y figuras de S1 a S13 sigue pendiente.}{S1 y S2 fuentes actuales; cotejo transversal S1–S13 pendiente.}

\observacion[tipo=acepta, sub=0]{09}{Títulos huérfanos al pie de página y cabeceras de tabla separadas del cuerpo.}{La condición de títulos huérfanos y continuidad de tablas solo se puede comprobar en la versión compilada; no hay PDF disponible.}{PDF compilado de la oferta; no disponible.}

\observacion[tipo=acepta, sub=0]{10}{La media firma del representante legal se superpone al pie de página y compite con la última línea de texto.}{La ubicación de la media firma y su relación con el pie de página no puede verificarse sin plantilla completa y PDF compilado.}{Plantilla completa y PDF compilado; no disponibles.}

\observacion[tipo=acepta, sub=0]{11}{Diagramas relegados a un capítulo final a media página, ninguno citado ni analizado desde el texto.}{S1 cita un organigrama; las imágenes de apoyo siguen en preparación y no se pueden cotejar en esta versión.}{S1 §1.2.1; archivo del organigrama y revisión visual pendientes.}

\observacion[tipo=acepta, sub=0]{12}{Se mencionan integrantes y organización interna del equipo en la oferta técnica en S13 y S3.}{Se mantiene la atribución por cargos corporativos. Los tres nombres ficticios de directivos fueron confirmados para conservarse en el ejercicio; S3 no los menciona y el S13 vigente no está disponible para cotejo.}{S1 §1.5; S3 vigente; S13 y T-8 final pendientes de cotejo.}

\observacion[tipo=acepta, sub=0]{13}{El Subdocumento 5 abre con la bitácora interna del grupo de trabajo.}{Verificado en la fuente S5 disponible: el documento abre con el modelo de datos y no con una bitácora interna.}{S5 §§5.1–5.2.}

\observacion[tipo=acepta, sub=0]{14}{El Subdocumento 2 cierra con una advertencia metodológica dirigida al profesor/evaluador.}{El cierre actual de S2 no contiene la advertencia metodológica dirigida al profesor/evaluador; cotejo textual favorable, falta revisión de la versión compilada.}{S2 §2.5; revisión de PDF pendiente.}

\observacion[tipo=acepta, sub=0]{15}{El folio de la propuesta final no es continuo y se reinicia en cada subdocumento.}{La numeración correlativa de páginas no se puede verificar: falta el archivo principal de compilación y no hay PDF final.}{Archivo principal de compilación y PDF final; no disponibles.}

\grupoObservaciones{Subdocumento 1, presentación de la empresa}

\observacion[tipo=acepta, sub=1]{16}{Una empresa de cuatro años y 22 personas promete presencia en cuatro nodos y reemplazo de hardware en menos de cuatro horas sobre 3.000 kilómetros sin indicar con quién}{Se declaran cuatro centros de soporte en el corredor de Ruta 5 y los convenios asociados, confirmados para S1. El plazo de sustitución menor a cuatro horas se presenta como compromiso de la oferta; su fuente y alcance se deben cotejar antes de atribuirlo como exigencia de las bases.}{1.6; cotejo de fuente y alcance del plazo pendiente}

\observacion[tipo=aclara, sub=1]{17}{Cita normativa errada: el D.S. 43 del MINSAL se invoca como norma de transporte cuando corresponde al D.S. 298}{Se distingue formalmente la aplicación copulativa de ambos reglamentos: el D.S. 298 rige el transporte terrestre en ruta para los 18 vehículos de carga peligrosa, mientras que el D.S. 43 rige el almacenamiento y permanencia en patios industriales y recintos de acopio donde operan dichas unidades.}{1.1.5}

\observacion[tipo=acepta, sub=1]{18}{Prosa de folleto publicitario: socio tecnológico de referencia en el Cono Sur y compromiso irrestricto}{Se purga la totalidad de expresiones promocionales y adjetivos subjetivos, sustentando la propuesta exclusivamente en especificaciones técnicas de ingeniería, proyectos comprobables y certificaciones auditables.}{1.1}

\observacion[tipo=acepta, sub=1]{19}{Líneas de negocio reducidas a un simple listado inconexo de tecnologías}{Se formalizan las tres líneas de servicio de la compañía (Telemetría de Borde, Plataformas Cloud Resilientes y Ciberseguridad Operacional), detallando el alcance de ingeniería, metodologías aplicadas y entregables verificables de cada una.}{1.1.2}

\observacion[tipo=acepta, sub=1]{20}{El modelo de gobierno corporativo es una lista de normas; faltan roles, comités, indicadores y ciclo de auditoría}{Se formaliza el modelo de gobierno institucional articulando tres comités colegiados permanentes (CCPI, CSIC, CEGCN), periodicidades de sesión, matriz de escalamiento técnico de cuatro niveles y ciclos anuales de auditoría bajo ISO 9001 e ISO/IEC 27001.}{1.3}

\observacion[tipo=acepta, sub=1]{21}{Ausencia de antecedentes financieros corporativos que acrediten solvencia para sostener 56 meses de contrato}{Los ratios y respaldos financieros se presentan en el Sobre N.º 1, conforme a lo confirmado para este ejercicio. S1 mantiene una remisión breve y no reproduce esos indicadores en la oferta técnica.}{S1 §1.4.3; cotejar la ubicación de los respaldos en el Sobre N.º 1 antes del cierre}

\observacion[tipo=acepta, sub=1]{22}{Los tres proyectos del Formulario T-6 no traen volúmenes comparables con 374 camiones y 96.000 viajes anuales}{Se acredita experiencia en tres contratos de gran envergadura (340, 310 y 390 unidades activas; entre 82.000 y 98.000 despachos anuales) con arquitecturas híbridas y dos proyectos con SLA mayor o igual a 99,5\%, estructurados en el Formulario T-6 independiente.}{1.4.1 y Formulario T-6}

\observacion[tipo=acepta, sub=1]{23}{La única figura del documento, el organigrama institucional, aparece reducida y no se explica en el texto}{Se desarrolla el texto explicativo sobre la estructura, los roles corporativos y la dotación confirmada. El archivo gráfico del organigrama está en preparación; al incorporarlo se verificará su legibilidad en la versión compilada.}{1.2; figura y revisión visual pendientes}

\grupoObservaciones{Subdocumento 2, comprensión del problema y de la necesidad}

\observacion[tipo=aclara, sub=2]{24}{No existe resumen ejecutivo de la propuesta; la sección 1.1 resume exclusivamente el problema}{S2 abre con un resumen que delimita el desafío, los antecedentes y el alcance del análisis del problema. El diseño de la solución se desarrolla en S3. Se debe cotejar contra T-7 y T-22 si esta separación satisface íntegramente la observación antes de cerrar.}{Resumen de apertura y 2.1; cotejo de T-7/T-22 pendiente}

\observacion[tipo=acepta, sub=2]{25}{Bloques 1.2.1 a 1.2.7 como títulos y tablas aisladas con prosa que repite adjetivos sin análisis de fondo}{Se sustituyen las tablas aisladas por un diagnóstico holístico continuo que jerarquiza impactos y dependencias causales en la cadena de valor, analizando subsidios cruzados, tacógrafos, descansos laborales y pasos cordilleranos.}{2.2}

\observacion[tipo=acepta, sub=2]{26}{Lenguaje inflado: riesgo de magnitudes incalculables, hemorragia financiera, desgobierno}{Se erradican adjetivos melodramáticos y prosa alarmista, reescribiendo el análisis en tono pericial corporativo, neutro, descriptivo y fundado en hechos fácticos comprobables.}{Subdocumento 2 completo}

\observacion[tipo=acepta, sub=2]{27}{Nota de blindaje económico dirigida al evaluador (pág. 11), indicio de uso de IA sin revisión}{Se elimina por completo la nota de advertencia económica dirigida al evaluador, salvaguardando la ficción institucional corporativa y los estándares de revisión humana.}{2.1 y 2.2}

\observacion[tipo=aclara, sub=2]{28}{No hay un solo dato investigado fuera del caso; la bibliografía los lista pero el texto no los utiliza}{S2 incorpora el marco de subcontratación, jornada, tránsito y tacógrafo, además del marco GLEC/ISO 14083, y describe consecuencias operativas o controles a confirmar. Las fuentes ya constan en la bibliografía compartida; la aplicabilidad jurídica debe revisarse antes del cierre.}{S2 §§2.2--2.3; referencias/referencias.bib; revisión jurídica pendiente}

\observacion[tipo=acepta, sub=2]{29}{La estacionalidad operativa se menciona pero no se cuantifica su efecto sobre jornada, sobreestadía ni cobertura}{S2 atribuye al caso la concentración estacional de demanda y las esperas de andén, pero aclara que el total anual no permite derivar un máximo diario de viajes de frío. Mantiene 72 horas como requisito de retención; 288 horas es una propuesta de D4, no una exigencia. La curva y su figura requieren cotejo de fuente.}{2.3; curva, figura y trazabilidad de parámetros pendientes}

\observacion[tipo=acepta, sub=2]{30}{El marco normativo aparece como citas sueltas sin explicitar la obligación de control que impone cada norma}{Se explicita la obligación de control técnico de cada norma: Ley 20.123 (responsabilidad solidaria), Art. 25 bis (descanso de 2 horas tras 5 horas continuas), D.S. 298 frente a D.S. 43 (SUSPEL), Ley No Chat 21.377 y Ley 21.719 de datos personales.}{2.2 y 2.3}

\observacion[tipo=acepta, sub=2]{31}{Faltan como actores el Fondo de Inversión (22\% de la propiedad), la autoridad laboral (DT) y la aseguradora de carga}{Se incorporan las fichas técnicas completas de los tres actores omitidos, analizando poder, interés, requerimiento legal y riesgo operacional asociado a cada uno en el ecosistema logístico.}{2.4 y Anexo 2.D}

\observacion[tipo=acepta, sub=2]{32}{Las tensiones de gobernanza se enumeran pero no se arbitran técnicamente}{Se arbitran formalmente las seis tensiones de gobernanza, sustentando la postura técnica y legal concluyente de audIT SpA en la normativa vigente (privacidad frente a fiscalización laboral bajo Ley 21.719).}{2.4.1}

\observacion[tipo=acepta, sub=2]{33}{Los 26 supuestos presentados son copia de las 26 decisiones del numeral 16.1 del pliego, no supuestos propios}{Se eliminan las copias de decisiones del pliego y se formulan supuestos auténticos de ingeniería de proyectos (SUP-01 a SUP-06) vinculados a la implementación del sistema.}{2.5.2}

\observacion[tipo=acepta, sub=2]{34}{No se declaran los supuestos propios de la propuesta técnica ni el impacto en caso de resultar falsos}{Se construye la Matriz de Supuestos de Ingeniería evaluando probabilidad, severidad, impacto operacional cuantitativo y plan de mitigación o contingencia técnica ante desviaciones.}{2.5.2}

\observacion[tipo=acepta, sub=2]{35}{El cierre de la sección 1.6.3 es una advertencia metodológica dirigida al profesor}{Se elimina la advertencia metatextual, sustituyéndola por la formalización de restricciones contractuales y exclusiones del diagnóstico.}{2.5.3 y Anexo 2.C}

\observacion[tipo=acepta, sub=2]{36}{Tono de paper académico y una sola figura en 32 páginas de documento}{Se revisa el tono y se agrega análisis técnico en el texto de S2. Las figuras previstas están en preparación y no se consideran incorporadas hasta recibir los archivos finales y verificar contenido, trazabilidad, legibilidad y ubicación en el PDF.}{2.1 a 2.4; entrega de figuras y revisión visual pendientes}

\grupoObservaciones{Subdocumento 3, esquema de solución y alcance}

\observacion[tipo=acepta, sub=3]{37}{No existe esquema de solución; cero figuras en 18 páginas de documento.}{Cotejo del checkout: S3 cita la Figura 3.1, pero el archivo gráfico referenciado no está disponible. Incorporación y revisión visual pendientes.}{S3 §3.3, Figura 3.1; recurso gráfico ausente.}

\observacion[tipo=acepta, sub=3]{38}{No hay modelo conceptual, componentes ni matriz de requerimiento a componente.}{S3 describe capas y contextos; no se encontró la matriz consolidada de los 42 requerimientos hacia componentes. Completar y adjuntar evidencia.}{S3 §3.3; matriz req-componente pendiente; T-17 aporta parte del mapeo.}

\observacion[tipo=acepta, sub=3]{39}{El capítulo abre directamente en el principio rector sin explicar el objetivo del subdocumento.}{Verificado parcialmente: S3 abre con objetivo, contexto y propósito. Falta el cotejo final de todas las referencias cruzadas.}{S3 §§3.1–3.2; cotejo de referencias internas pendiente.}

\observacion[tipo=acepta, sub=3]{40}{La sección 1.4 remite a una Sección 4 que no existe en el documento.}{El archivo actual usa numeración 3.x y no la sección 1.4 observada; queda pendiente revisar todas las referencias cruzadas contra la versión que se entregará.}{S3 numeración 3.x; revisión de referencias internas pendiente.}

\observacion[tipo=acepta, sub=3]{41}{La Etapa 1 es una tabla de 9 capacidades sin identificador y la Etapa 2 un párrafo, sin dependencias ni hitos externos.}{La Tabla 3.1 divide capacidades entre etapas; no se encontraron identificadores unívocos, dependencias, horizonte 2029 ni análisis de capacidad de absorción. Pendiente completar o retirar esas afirmaciones.}{S3 §3.2.1, Tabla 3.1; completar IDs, dependencias y horizonte.}

\observacion[tipo=acepta, sub=3]{42}{Exclusiones copiadas, supuestos sin listado formal y restricciones no recogidas.}{S3 §3.2.2 formaliza exclusiones, supuestos y restricciones, pero presenta 288 h como exigencia. El requisito confirmado es 72 h (RT-03.10); 288 h es una ampliación propuesta por D4. Corregir antes del cierre.}{S3 §3.2.2; requisito 72 h RT-03.10 confirmado; corregir referencia a 288 h.}

\observacion[tipo=acepta, sub=3]{43}{La decisión sobre el sistema de gestión legacy de 2013 no se menciona ni una vez en el capítulo.}{S3 menciona el ERP de 2013 y su integración, pero no se encontró allí el ADR-01 ni una decisión formal Strangler Fig/CDC. Pendiente incorporar evidencia o ajustar la respuesta.}{S3 §3.3 y S4 §4.1.5; ADR-01/decisión formal no encontrada.}

\observacion[tipo=acepta, sub=3]{44}{Sin visión de arquitectura empresarial ni análisis de interoperabilidad, escalabilidad, seguridad o performance.}{S3 expone capas y tecnologías; falta cotejar y documentar por separado la interoperabilidad, escalabilidad, seguridad y rendimiento que afirma esta respuesta.}{S3 §§3.3–3.4; análisis trazable pendiente.}

\observacion[tipo=acepta, sub=3]{45}{Se declaran 42 requisitos, se muestran solo 12 y el Formulario T-12 no fue entregado.}{La fuente T-12 disponible contiene 28 RF y 14 RNF. No está el PDF final en este checkout; exportar y cotejar la versión de entrega.}{AUDIT-Formulario-T-12.md, RF-001–028 y RNF-001–014; PDF pendiente.}

\observacion[tipo=acepta, sub=3]{46}{Los 12 requisitos mostrados no tienen actor, precondición, resultado esperado ni prioridad (numeral 17.1 del Caso).}{La fuente T-12 contiene 42 requisitos, pero no se ha verificado que cada ficha cumpla los 12 campos obligatorios; el esquema RNF usa campos distintos. Revisar y entregar el formulario final.}{AUDIT-Formulario-T-12.md; cotejo de los campos por ficha y PDF pendiente.}

\observacion[tipo=acepta, sub=3]{47}{No hay registro de reglas de negocio: cómputo de jornada, excepciones, tiempo libre de espera ni retornos vacíos.}{No se encontró en S3 un registro formal RN-01 a RN-12. Incorporar las reglas de negocio y su modelado, o ajustar la respuesta.}{S3; registro RN-01–RN-12 no encontrado.}

\observacion[tipo=acepta, sub=3]{48}{No hay matriz de trazabilidad de origen a requerimiento a componente a prueba.}{T-17 contiene mapeos RF→componente→pruebas; falta cotejar la cadena completa origen contractual→T-12→componente→prueba en una matriz consolidada.}{Formulario T-17 §3; completar traza desde origen contractual.}

\observacion[tipo=aclara, sub=3]{49}{Implementación, implantación y operación no existen (nada sobre CI/CD, migración, cutover, RTO/RPO ni DR).}{S3 no desarrolla CI/CD, migración, cutover, RTO/RPO ni DR. S6 está en el checkout; S7 y S8 no están disponibles aquí. Cotejar T-7/T-22 y adjuntar los documentos que sustenten la aclaración.}{S3; S6 disponible; S7/S8 ausentes del checkout; verificar T-7/T-22.}

\observacion[tipo=acepta, sub=3]{50}{Para camiones con GPS ajeno se promete jornada acreditada con equipos sin identificación; la evidencia es mera atestación.}{S3 describe evidencia sellada y niveles de evidencia, y T-17 menciona atestación externa; no se encontró la cascada completa por GPS propio, GPS de terceros y ausencia de GPS. Pendiente especificarla en S3/S5.}{S3 §3.4.1, T-17 RF-003 y S5; cascada probatoria pendiente.}

\observacion[tipo=acepta, sub=3]{51}{Se afirma que 29 criterios están comprometidos y solo ocho tienen meta, hito y medición.}{S3 enumera 29 criterios y desarrolla solo categorías y ejemplos seleccionados; no se encontró la matriz completa con meta, hito y método para CA-01 a CA-29. Pendiente.}{S3 §3.2.4; matriz CA-01–CA-29 con metas, hitos y métodos pendiente.}

\observacion[tipo=acepta, sub=3]{52}{Códigos sin definir en su primera aparición (decisión D-02 y supuesto S-09).}{No se encontró en S3 un glosario que defina D-XX, SUP-XX y RN-XX en su primera aparición. Pendiente incorporar y verificar.}{S3; glosario D-XX/SUP-XX/RN-XX no encontrado.}

\grupoObservaciones{Subdocumento 4, arquitectura lógica}

\observacion[tipo=acepta, sub=4]{53}{El diagrama de las ocho capas no muestra interfaces y no se explica qué contiene cada capa ni cómo se conectan}{Dos tablas declaran el contenido y los componentes de cada capa, y el texto explica que una petición atraviesa las seis capas de la pila sin acceso directo a la base de datos. Se incorpora: diagrama redibujado con las interfaces entre capas.}{4.1.3}

\observacion[tipo=acepta, sub=4]{54}{Componentes que aparecen solo como ícono, terminal de torre y pantalla de terreno, sin decir qué equipo es, qué software ejecuta, quién lo usa ni con qué se conecta}{La capa de presentación declara el portal, la aplicación en los cuatro perfiles de RT-17.01 del Caso y las vistas de portería y de taller. Se incorpora: explicación de cada componente, dónde corre y con qué se conecta.}{4.1.3}

\observacion[tipo=acepta, sub=4]{55}{No hay matriz de requerimiento a componente, ningún requerimiento funcional del Subdocumento 3 se ubica en un componente}{La trazabilidad cruza cada componente crítico entre la capa lógica, el nodo físico y el Formulario T-11. Se incorpora: matriz de requerimiento funcional del T-12 a componente.}{4.1 y 4.2.3}

\observacion[tipo=acepta, sub=4]{56}{No se mencionan los principios SOLID, las arquitecturas de referencia se reducen a ISO 42010 y los ambientes del ciclo de desarrollo están ausentes}{Los cinco ambientes, de desarrollo a recuperación, van cada uno en su suscripción y se levantan desde el mismo código de infraestructura. Se incorpora: principios SOLID y arquitecturas de referencia.}{4.1 y 4.2.6}

\observacion[tipo=acepta, sub=4]{57}{El diagrama de secuencia de la asignación bloqueante está huérfano, el texto no lo cita ni lo recorre}{El texto recorre la asignación desde la puerta de enlace hasta la persistencia, y la reconexión tras 72 horas se dimensiona tramo por tramo. Se incorpora: secuencias de emisión sin cobertura y de liquidación.}{4.1.9 y 4.2.8}

\observacion[tipo=acepta, sub=4]{58}{Frameworks y tecnologías viven en las figuras, el texto no las decide ni las justifica}{El catálogo de nube declara nivel, configuración y subred de cada servicio. Se incorpora: tabla de lenguajes, frameworks y middleware con versión, fin de soporte y alternativas descartadas.}{4.1.1 y 4.2.4}

\observacion[tipo=acepta, sub=4]{59}{El modelo táctico no trae límites de contexto ni eventos, aunque la sección lo promete}{Se incorpora: límites de contexto y eventos del modelo táctico.}{4.1.16}

\observacion[tipo=acepta, sub=4]{60}{El estrangulamiento del sistema de 2013 no tiene registro de decisión ni compara alternativas ni fundamento económico}{Se incorpora: registro de decisión con reemplazo completo, conservar e integrar y paquete de mercado, con su fundamento económico.}{4.1.6}

\observacion[tipo=acepta, sub=4]{61}{El descarte de la banda ancha satelital queda registrado pero no aparece como registro de decisión}{El equipo a bordo usa mensajería satelital Iridium de hasta 340 bytes en los tramos sin señal. Se incorpora: registro de decisión del descarte de la banda ancha satelital.}{4.1.6 y 4.2.2}

\observacion[tipo=acepta, sub=4]{62}{Se afirma que el estilo arquitectónico se justifica con la volumetría sin justificarlo, y el numeral 2.3 de las Transversales pide comparar estilos}{Se incorpora: comparación de estilos del numeral 2.3 transversal.}{4.1 y 4.1.1}

\observacion[tipo=acepta, sub=4]{63}{El stack tecnológico no ha sido validado}{La sección 4.1.1 confirma el bus y el motor de series que usa la arquitectura física y declara el retiro de Azure Cache for Redis en 2028. Se incorpora: versión, fin de soporte y prueba de concepto de cada producto.}{4.1.1}

\grupoObservaciones{Subdocumento 4, arquitectura física}

\observacion[tipo=acepta, sub=4]{64}{Falta el diagrama de arquitectura física, la figura son íconos y cajas sin nivel de servicio, redes virtuales, subredes, nodos, enlaces ni anchos de banda}{Un diagrama general muestra los tres planos con sus enlaces y capacidades, y otro la región primaria con sus redes, subredes, servicios con su nivel y conexiones}{4.2 y 4.2.5}

\observacion[tipo=acepta, sub=4]{65}{No se dice qué componentes de Azure se usan ni con qué configuración, y las figuras se contradicen entre motores de series y entre buses}{El catálogo declara cada servicio con su nivel, configuración y subred. Las series van en PostgreSQL con TimescaleDB, en un servidor propio, y el bus es Event Hubs Premium con protocolo Kafka}{4.2.4}

\observacion[tipo=acepta, sub=4]{66}{La sala de 26 m² se compara ítem por ítem pero falta plano, carga eléctrica y térmica}{Se entregan el plano con las zonas del RT-06.03, el balance eléctrico y térmico con factor de potencia y PUE, y el dimensionamiento de las UPS y del grupo electrógeno}{4.3.1}

\observacion[tipo=acepta, sub=4]{67}{Ningún diagrama se cita desde el texto}{Cada figura se cita antes de aparecer y se explica después, en la sección que la usa}{Subdocumento 4 completo}

\observacion[tipo=acepta, sub=4]{68}{La sección promete declarar cada producto con versión, fin de soporte y plan de actualización y no declara ninguno}{El Formulario T-11 declara el ciclo de vida de cada equipo y servicio. Se incorpora: versión y fin de soporte del software.}{4.1.1 y Formulario T-11}

\observacion[tipo=acepta, sub=4]{69}{El Anexo A no es el T-11, y el Capítulo 11 del Caso exige qué comprar, cuánto y con qué características}{El Formulario T-11 va en archivo propio, con componente, producto, características, ubicación, cantidad, quién lo adquiere, ciclo de vida y justificación del emplazamiento}{Formulario T-11}

\observacion[tipo=acepta, sub=4]{70}{Sin cantidades, el parque oscila entre 374 y «según adhesión», y nunca se explicita que los equipos de terceros no se reemplazan}{Las cantidades salen de las poblaciones del Caso: 182 equipos a bordo más 19 de reposición, 572 tarjetas y 210 balizas más 21. Los 192 equipos de terceros no se reemplazan, y el crecimiento a tres años va en filas aparte}{4.2.1 y Formulario T-11}

\observacion[tipo=acepta, sub=4]{71}{Se dice que hay dieciocho componentes en la región primaria y la tabla suma veinte, y la matriz cubre 14 de los 34}{El conteo sale del Formulario T-11: 38 componentes, 23 de ellos en la región primaria, y la matriz de trazabilidad cruza los componentes críticos}{4.2.3}

\observacion[tipo=acepta, sub=4]{72}{La región secundaria no tiene nombre y el Artículo 16.3 exige declararla}{La región secundaria es Azure Brazil South, en espera, a unos 2.586 km de la primaria, con el mecanismo de réplica de cada servicio}{4.3.2}

\observacion[tipo=acepta, sub=4]{73}{Los 8 GB del dispositivo no tienen cálculo, y las 288 horas de cierre fronterizo dan cerca de 40 MB}{La capacidad se deriva por componente. Se requieren 2.374 MB, el búfer de 288 horas ocupa 38,4 MB y los 8 GB son la menor configuración del equipo}{4.2.2}

\observacion[tipo=acepta, sub=4]{74}{La sincronización en 20 minutos se repite como requisito y no se dimensiona, y cientos de unidades saliendo de una sombra no caben en esa ventana}{Trescientos camiones que salen a la vez de una sombra sincronizan 57.600 mensajes en 9,6 minutos con IoT Hub S1 de dos unidades, dentro de los 20 del RT-03.13}{4.2.8}

\observacion[tipo=acepta, sub=4]{75}{Faltan los eventos por segundo, el TPS del peak de asignación, el volumen transaccional anual, los datos a migrar, el ancho de banda por terminal y la mesa de ayuda}{Se declaran los eventos por segundo, las transacciones de la asignación, el volumen transaccional anual con la evidencia almacenada y el ancho de banda de terminales y de San Bernardo. La mesa de ayuda se dimensiona en el Subdocumento 10. Se incorpora: volumen de los datos históricos a migrar.}{4.2.8}

\observacion[tipo=acepta, sub=4]{76}{La restricción 8 fija al sistema contable como único emisor del documento electrónico y el mecanismo propuesto convierte al dispositivo en emisor}{El sistema contable queda como único emisor y ningún equipo del Formulario T-11 emite documentos tributarios. Se incorpora: mecanismo de emisión sin cobertura. Se incorpora: aclaración sobre la numeración de la consulta citada.}{4.1 y 4.2.2}

\observacion[tipo=acepta, sub=4]{77}{El equipo de TI de nueve personas no se menciona y el RT-03.05 se cita sin listar los servicios administrados}{El catálogo lista los servicios administrados y los justifica con el tamaño del área de tecnología del mandante, nueve personas para cinco terminales, dos talleres y la flota}{4.2.4}

\observacion[tipo=acepta, sub=4]{78}{Se tratan los tres proveedores de GPS y los 34 camiones sin equipo, pero no se dice qué se instala en esos 34 ni cómo conviven los dispositivos propios con los ajenos}{Los 34 camiones de terceros sin equipo llevan el equipo audIT, igual que los 148 propios. Los 192 de terceros con equipo lo conservan, sus datos llegan por las plataformas de sus proveedores y se fija el estándar mínimo para homologarlos}{4.2 y 4.2.1}

\observacion[tipo=acepta, sub=4]{79}{La recuperación ante desastres y la continuidad se separan con argumento sísmico, pero no se fundamenta su factibilidad}{La recuperación usa Brazil South con réplica declarada por servicio, tres defensas para la telemetría y conmutación en cuatro pasos. La continuidad sin enlace descansa en San Bernardo, los terminales y los camiones, con cada función declarada y su procedimiento supletorio}{4.3.2}

\observacion[tipo=acepta, sub=4]{80}{Los respaldos aparecen en una sola fila y no hay modelo Zero Trust explícito ni superficie de exposición}{Los respaldos siguen el esquema 3-2-1-1-0 con copia inmutable fuera de sitio, y la segmentación deja las subredes de aplicación y datos sin exposición a Internet. Se incorpora: modelo Zero Trust y superficie expuesta.}{4.1, 4.2.5 y 4.2.6}

\grupoObservaciones{Subdocumento 5, modelo y gestión de datos}

\observacion[tipo=acepta, sub=5]{81}{El diccionario cubre 6 entidades y ninguna es la evidencia de jornada, exigida para resistir impugnaciones.}{S5 define DOCUMENTO_RESPALDO con hash SHA-256 y almacenamiento WORM, pero no modela la entidad EvidenciaJornada ni los campos probatorios descritos en esta respuesta. Pendiente completar el diccionario.}{S5 §5.1.4, Tabla 5.5; entidad EvidenciaJornada pendiente.}

\observacion[tipo=acepta, sub=5]{82}{No hay diagrama Entidad-Relación (DER) en el documento.}{S5 referencia la Figura 5.1 como modelo conceptual y relacional, pero el archivo gráfico no está disponible en el checkout. No se puede verificar que sea un DER completo; imagen y revisión pendientes.}{S5 §5.1, Figura 5.1; recurso gráfico ausente.}

\observacion[tipo=acepta, sub=5]{83}{Contradicción de residencia: copia fuera de sitio en Azure East US 2 mientras S4 sostiene residencia en Chile.}{S5 no declara una región cloud ni contiene la sección 5.9 citada. S4 declara Chile Central y Brazil South; alinear S5 con S4 antes de afirmar que se corrigió la residencia.}{S4 §4.3.1 y §4.3.2; S5 sin declaración de residencia.}

\observacion[tipo=acepta, sub=5]{84}{El motor de bases de datos de series de tiempo y el bus de mensajería quedan indefinidos.}{S5 menciona TimescaleDB/Kafka, pero su Tabla 5.8 menciona TimescaleDB/Azure Event Hubs; S3 también nombra Event Hubs. Resolver y documentar el bus de mensajería elegido.}{S5 Tabla 5.1 y Tabla 5.8; resolver Kafka frente a Azure Event Hubs.}

\observacion[tipo=acepta, sub=5]{85}{La sección 1.1 es la bitácora interna del grupo de trabajo.}{Verificado: S5 comienza con la descripción conceptual del modelo de datos y no con una bitácora interna.}{S5 §5.1.}

\observacion[tipo=acepta, sub=5]{86}{Prosa hinchada y escasez de figuras (dos figuras en 21 páginas), no citadas desde el texto.}{S5 referencia las Figuras 5.1 a 5.6, pero los seis archivos gráficos no están disponibles en el checkout. La ubicación, legibilidad e interpretación visual quedan pendientes de cotejo.}{S5 Figuras 5.1–5.6; recursos gráficos y revisión visual pendientes.}

\observacion[tipo=acepta, sub=5]{87}{El tratamiento de revocación del consentimiento frente a la retención legal queda pendiente de validar.}{S5 incluye fecha de revocación y una matriz de retención de cinco años, pero no contiene la sección 5.11 citada ni el análisis jurídico expreso que afirma la respuesta. Completar el fundamento y corregir la referencia.}{S5 Tablas 5.6 y 5.10; sección 5.11 y fundamento jurídico pendientes.}

\grupoObservaciones{Subdocumento 13, innovaciones}

\observacion[tipo=acepta, sub=13]{88}{La innovación tipo 3 es el RT-03.10, el RT-03.13, los criterios 8 y 9 y la decisión 4, una funcionalidad exigida presentada como innovación, que el Artículo 30 excluye}{La innovación 3 pasa a ser el semirremolque conectado, que ninguna exigencia de las bases pide. La operación sin cobertura queda como alcance base en la sección 4.2.2 del Subdocumento 4}{13.3}

\observacion[tipo=acepta, sub=13]{89}{La innovación tipo 2 es el método que imponen las restricciones del Capítulo 6 y el numeral 13.3}{Se incorpora: reformulación o reemplazo de la innovación 2.}{13.2}

\observacion[tipo=acepta, sub=13]{90}{Las fichas T-19 con los siete elementos del Artículo 29 se difieren a la propuesta final, aunque el Informe 1 las incluye}{Las cinco fichas van en el Formulario T-19, en archivo propio}{Formulario T-19}

\observacion[tipo=acepta, sub=13]{91}{Los tipos 2 y 3 no traen madurez, impacto económico ni riesgo de adopción, el tipo 5 no trae riesgo y el tipo 1 solo declara inversión incremental}{Cada sección se ordena por los siete elementos del Artículo 29, y la innovación 3 los declara completos. Se incorpora: madurez, impacto y riesgo de la innovación 2. Se incorpora: costo operacional y contingencia de la innovación 1.}{13.1 a 13.5 y Formulario T-19}

\observacion[tipo=acepta, sub=13]{92}{En la tabla del tipo 3 las líneas base están inventadas y la meta está redactada al revés respecto del RT-03.13}{La nueva innovación 3 declara como línea base lo que el sistema actual no registra, sin cifras ajenas al Caso, y mide sus metas con datos de la propia solución}{13.3}

\observacion[tipo=acepta, sub=13]{93}{El tipo 5 usa cifras que nadie midió, la línea base NASA-TLX, el 22 \% de adherencia, el 1,2 \% de ahorro y el costo por evento}{Se retiran las cuatro cifras. La carga NASA-TLX y la adherencia pasan a línea base medida en la Etapa 1, y el efecto en combustible se mide en la misma etapa}{13.5}

\observacion[tipo=acepta, sub=13]{94}{Las 280 horas de desarrollo declaradas son información de esfuerzo dentro de la Oferta Técnica, contra el Artículo 50.2}{Se retiran}{13.5}

\observacion[tipo=acepta, sub=13]{95}{El tipo 4 cuantifica un recupero anual en pesos}{Se retiran los montos. El recupero se valoriza en la Oferta Económica y no se compromete como ingreso}{13.4}

\observacion[tipo=acepta, sub=13]{96}{El producto propio audIT EdgeHub se invoca sin definirlo en ningún subdocumento}{Se retira el nombre. La innovación 5 remite al equipo a bordo que especifica la sección 4.2.2 del Subdocumento 4}{13.5}

\observacion[tipo=acepta, sub=13]{97}{No hay trazabilidad con el flujo de caja para ninguna innovación, y la tipo 4 debía reflejarse en la estructura de costos}{La innovación 3 se refleja como inversión en equipamiento en la Etapa 1 y como costo de reposición en la operación. Se incorpora: partidas de las innovaciones 1 y 4. Se incorpora: partida de la innovación 2. Se incorpora: partida de la innovación 5.}{13.1 a 13.5}

\observacion[tipo=aclara, sub=13]{98}{Declarar investigación adicional requerida en cada ficha equivale a declarar que la innovación no está diseñada}{Las bases ordenan que «si alguna requiere investigación adicional, debe declararse» (\form{T-22}{68}), y audIT usó esa figura sin reemplazar el diseño. En este informe cada pendiente queda dentro del elemento del Artículo 29 al que pertenece}{13.1 a 13.5}

\observacion[tipo=acepta, sub=13]{99}{El punto 1.6 sobre estado de elaboración y responsables es organización interna, ininteligible para el mandante}{Se eliminan la sección y los nombres}{Subdocumento 13 completo}

\observacion[tipo=acepta, sub=13]{100}{Cita al autor de las bases como fuente bibliográfica, tabla partida entre páginas y cero figuras}{Las bases se citan por documento, numeral y página. El capítulo trae dos figuras citadas antes y explicadas después}{Subdocumento 13 completo}

\end{tablaObservaciones}

\marcador{Las 100 observaciones ya están incluidas. Cotejar cada respuesta, estado y referencia contra la matriz fuente y la evidencia oficial antes de cerrar la tabla del Artículo 46; las respuestas incorporadas desde el borrador no se consideran verificadas.}
